\chapter{The synthetic fleet: SynCFM56}
\label{ch:fleet}

This chapter documents phase F2: the generator that turns the F1 twin into a 100-engine
run-to-failure fleet with complete ground truth, the audits that gate it (one of which failed
on first pass and forced a design change --- documented in full in \cref{sec:audits}), and the
declaration of milestone H2.

\section{Generator architecture: linearize, then audit}
\label{sec:fleet-arch}

A fleet of 100 engines flying $\sim$16\,000 cycles each, two snapshots per flight, would
require over three million full pyCycle solves per fleet realization --- unusable for
iteration. Instead, the generator exploits exactly what F1 built: every snapshot is produced by
the linearization
\begin{equation}
  z(u, \mathbf{x}) \;=\; z_0(u)\,\bigl(1 + \mathbf{H}(u)\,\mathbf{x}/100\bigr),
  \label{eq:generator}
\end{equation}
where $z_0(u)$ is the healthy baseline and $\mathbf{H}(u)$ the influence coefficient matrix,
both linearly interpolated between the ICM grid points that bracket the sampled operating
condition (takeoff snapshots: in $\Delta T_{\mathrm{ISA}}$; cruise snapshots: in commanded N1).
N1 itself carries zero deviation by construction --- it is the held power-setting parameter,
not a health channel. The whole fleet generates in 37 seconds on all cores (one worker process
per engine, norm N1), with per-engine deterministic random seeds so the output is reproducible
regardless of process scheduling.

Two properties of this choice must be kept distinct. First, the dataset is
\emph{self-consistent by construction}: the recorded ground truth $\mathbf{x}$ is exactly the
state that generated the measurements, so estimator evaluations against truth are exact
regardless of any linearization error. Second, \emph{fidelity to the full nonlinear physics} is
an empirical question --- which is why the first gate audit re-solves full pyCycle at states
sampled from the generated fleet (\cref{sec:audit-nonlin}).

\section{Ground-truth trajectories}
\label{sec:trajectories}

Each engine draws a degradation ``personality'' from the fault catalog
(\code{conf/fault_catalog.yaml}, the versioned contract): a log-normal severity multiplier per
mechanism, a wash schedule, Poisson-arriving foreign-object damage, optional sensor drifts, and
---in roughly half the fleet--- one \emph{acute fault episode} (\cref{sec:audit-difficulty}
explains why these exist). The mechanisms follow the physical profiles of
\cref{tab:degradation}: saturating-exponential fouling, slow linear erosion, bilinear
break-in of tip clearances, linearly accelerating hot-section deterioration.

Two implementation details worth recording:
\begin{itemize}
  \item \textbf{Wash sawtooth, exact.} Between washes the fouling level relaxes exponentially
    toward its asymptote; each wash multiplies the current level by $(1-r)$ with recovery $r$
    drawn per event. The trajectory is computed segment-wise with the exact exponential
    solution. A tempting shortcut --- permanently subtracting the removed amount --- was
    implemented first and rejected: it makes the trajectory approach
    $\mathrm{asymptote} - \mathrm{removed}$ and under-predicts late-life fouling.
  \item \textbf{Composability of truth.} The generator stores not only the total
    $\mathbf{x}(n)$ but each mechanism's separate contribution, so any later analysis can ask
    not just ``how degraded'' but ``degraded by \emph{what}''.
\end{itemize}

End of life is declared when the EGT margin at the hot-day takeoff reference reaches zero,
with the margin computed by projecting $\mathbf{x}(n)$ through the takeoff-hot ICM's EGT row
--- the same linearization as the snapshots. New-engine margins are drawn per engine
($85\pm6$\,\si{\degreeCelsius}). \Cref{fig:egtm} shows the resulting fleet behavior: the
canonical fast-early / slow-late erosion with wash sawtooth, and lives spanning
10\,500--20\,700 cycles.

\begin{figure}[htbp]
  \centering
  \includegraphics[width=0.85\textwidth]{egtm_sawtooth}
  \caption{EGT-margin trajectories of five engines. Wash events produce the sawtooth
  recoveries; engine-to-engine spread comes from the hierarchical severity multipliers and the
  drawn new-engine margins. End of life is the zero crossing. Auto-generated from the fleet.}
  \label{fig:egtm}
\end{figure}

\begin{figure}[htbp]
  \centering
  \includegraphics[width=0.85\textwidth]{health_trajectory}
  \caption{Ground-truth health parameters of one engine with an acute fault episode (dashed
  marker). Vertical gray lines are washes (visible as sawtooth in the HPC traces); the
  hot-section traces show the monotone efficiency loss and flow-capacity opening.
  Auto-generated.}
  \label{fig:trajectory}
\end{figure}

\section{Snapshots and the measurement model}
\label{sec:snapshot-model}

Each flight cycle emits two reports, mimicking airline ACARS practice: a takeoff snapshot (sea
level, ambient offset drawn from a climate distribution) and a cruise snapshot
(FL350/M0.78, commanded N1 drawn per flight). Truth channels pass through the sensor model of
the catalog: per-channel Gaussian noise (absolute or percentage $\sigma$), quantization to the
recording resolution, additive drift ramps on the affected engines (37 of 100 --- matching the
catalog's per-sensor probabilities summing to 0.36), and missing data as isolated losses plus
ACARS outage bursts (4.1\,\% of snapshots). Every row carries, alongside the measured values:
the true channel values, the full ground-truth health vector, the EGT margin, the remaining
useful life, the isolation label, and drift flags. \Cref{tab:fleet} summarizes the frozen
dataset.

\begin{table}[htbp]
  \centering\small
  \caption{SynCFM56 v1.0 at a glance. Auto-generated from the fleet artifacts.}
  \label{tab:fleet}
  % AUTOGENERATED by scripts/make_report_assets.py — do NOT edit by hand
\begin{tabular}{lr}
\toprule
Property & Value \\
\midrule
  Engines (train / val / test) & 100 (70 / 10 / 20) \\
  Snapshot rows (takeoff + cruise per flight) & 1\,584\,808 \\
  Life [cycles]: median (min--max) & 15442 (10515--20738) \\
  Censored engines & 0 \\
  Engines with an acute fault episode & 49 (all 6 target parameters) \\
  Engines with a sensor drift & 37 \\
  Wash EGTM recovery: mean, frac.\ positive & 2.8 °C, 99.8\,\% \\
  End-of-life $\Delta$EGT / $\Delta W_f$ (all engines correct sign) & $+68$ K / $+5.8$\,\% \\
  Measured-noise residual vs.\ catalog (EGT) & 2.51 K vs.\ 2.5 K \\
\bottomrule
\end{tabular}

\end{table}

\section{Gate audits --- including the one that failed}
\label{sec:audits}

\subsection{Nonlinearity of the generator}
\label{sec:audit-nonlin}

Eighty health states were sampled from the generated fleet (stratified by life fraction,
half from the last 30\,\% of life where deviations are largest) and re-solved with full
pyCycle at the exact ICM grid conditions, so the comparison isolates nonlinearity from
interpolation. \Cref{tab:nonlin} reports the per-channel error of \cref{eq:generator} against
the full solution, next to the sensor noise it must be judged against: medians of
0.01--0.33\,\%, tails (P95) at or below 1\,\%, concentrated at end of life. Combined with the
self-consistency property of \cref{sec:fleet-arch}, the linearization is adequate: the residual
infidelity to full physics is at the noise floor, identical for both EHM pipelines, and
irrelevant to truth-referenced scoring.

\begin{table}[htbp]
  \centering\small
  \caption{Nonlinearity audit: linear generator versus full pyCycle solves at fleet-sampled
  health states. Auto-generated.}
  \label{tab:nonlin}
  % AUTOGENERATED by scripts/make_report_assets.py — do NOT edit by hand
\begin{tabular}{llrrrr}
\toprule
Condition & Channel & median [\%] & P95 [\%] & max [\%] & noise $\sigma$ [\%] \\
\midrule
  Cruise & N2 & 0.015 & 0.036 & 0.039 & 0.05 \\
  Cruise & $W_f$ & 0.238 & 0.592 & 0.760 & 0.50 \\
  Cruise & EGT & 0.336 & 0.963 & 1.041 & 0.21 \\
  Takeoff, hot day & N2 & 0.038 & 0.091 & 0.095 & 0.05 \\
  Takeoff, hot day & $W_f$ & 0.135 & 0.449 & 0.483 & 0.50 \\
  Takeoff, hot day & EGT & 0.245 & 0.860 & 0.933 & 0.18 \\
\bottomrule
\end{tabular}

\end{table}

\subsection{Difficulty: the audit that failed, and what it taught}
\label{sec:audit-difficulty}

The difficulty gate is designed to prevent a trivially easy dataset from flattering the AI: a
deliberately weak classifier (multinomial logistic regression on the 16 raw measured channels)
must \emph{not} solve the isolation task; the pre-set threshold is 60\,\% accuracy.

\textbf{First pass: failed.} With the original labeling --- the dominant \emph{chronic}
mechanism per snapshot --- the trivial classifier scored \textbf{81.6\,\%} (majority-class
baseline 54.6\,\%). Root-cause analysis: chronic mechanisms co-evolve with age in \emph{every}
engine (fouling saturates early, hot-section deterioration accumulates late), so the chronic
label is largely an \emph{age proxy}, and raw absolute channel levels encode age directly. The
task the label defined was not fault isolation at all. Two responses were possible: relax the
threshold (cosmetic) or fix the design (structural). The structural fix was taken.

\textbf{Design change: acute fault episodes.} What hypothesis H2 actually tests is the
isolation of \emph{localized faults with overlapping signatures} --- the regime where the ICM
observability analysis proved cockpit sensors nearly blind (minimum signature angle
1.3$^{\circ}$, \cref{tab:confusable}). The catalog therefore gained discrete
\emph{single-parameter} fault episodes: a fast ramp (50--500 cycles) to a sustained magnitude
on one health parameter, drawn from six targets that deliberately include the confusable pairs
($\eta_{\mathrm{HPC}}$ vs.\ $\eta_{\mathrm{HPT}}$, $\eta_{\mathrm{HPT}}$ vs.\
$\Gamma_{\mathrm{HPT}}$), injected in $\sim$50\,\% of engines at a random point of their lives.
The gated task became: classify each snapshot as \emph{no acute fault} or by the faulted
parameter --- a 7-class isolation problem.

\textbf{The audit also caught a generator bug.} The first acute implementation drew the episode
onset as a fraction of the \emph{maximum simulation horizon} (30\,000 cycles) rather than of
the engine's actual life ($\sim$16\,000): most episodes landed beyond end of life and only
15/100 engines showed one. The fix computes a chronic-only provisional life first, then places
the onset as a fraction of it (two-pass generation); the final fleet has 49/100 acute engines
with all six target parameters represented.

\textbf{Second pass: passed.} On the corrected fleet the trivial classifier achieves
\textbf{54.0\,\%} accuracy (macro-F1 0.30) on the gated isolation task --- under the 60\,\%
threshold, exactly as the signature-angle physics predicts (\cref{tab:difficulty-audit}). The
chronic-mechanism label is retained in the dataset and reported as a secondary,
explicitly-not-gated metric (81.6\,\% accuracy --- high, expected, and uninformative about
isolation ability). The episode is recorded here in full because it is evidence the gates are
real: a pre-registered criterion failed, the failure was diagnosed to a design flaw rather than
argued away, and the fix strengthened the benchmark.

\begin{table}[htbp]
  \centering\small
  \caption{Difficulty audit on the final fleet (trivial classifier: multinomial logistic on
  raw measured channels). Auto-generated.}
  \label{tab:difficulty-audit}
  % AUTOGENERATED by scripts/make_report_assets.py — do NOT edit by hand
\begin{tabular}{lrrrl}
\toprule
Trivial-classifier task & Accuracy [\%] & macro-F1 & Gate [\%] & Verdict \\
\midrule
  Acute-fault isolation (gated, 7 classes) & 54.0 & 0.299 & $<60$ & PASS \\
  Chronic-mechanism label (secondary, age proxy) & 80.1 & 0.607 & --- & reported \\
\bottomrule
\end{tabular}

\end{table}

\subsection{Realism}
\label{sec:audit-realism}

Quantitative checks of macroscopic behavior, all in \cref{tab:fleet}: the wash sawtooth is
visible (mean EGTM recovery $+2.8$\,\si{\degreeCelsius}, positive at 99.75\,\% of the 400
sampled wash events); end-of-life deviations carry the physical deterioration signature in
100\,\% of engines ($\Delta$EGT $+68$\,K, $\Delta W_f$ $+5.7$\,\%); the measured-noise
residual on drift-free engines reproduces the catalog $\sigma$ (2.51 vs.\ 2.50\,K); lives
center on 15\,400 cycles with a realistic severe-degrader tail and no censored engines.

\section{Milestone H2: declared closed}
\label{sec:h2}

All gate criteria are met: realism and difficulty audits passed (the latter after the
documented design change), nonlinearity quantified at the noise floor, split-by-engine
leakage checks automated in CI, and the dataset datasheet written
(\code{docs/syncfm56-datasheet.md}). SynCFM56 v1.0 is frozen: regeneration is deterministic
from the catalog and seed, so the dataset is fully specified by the repository state. Known
limitations, declared for downstream chapters: linearized generation (audited,
\cref{tab:nonlin}); two fixed snapshot conditions with ambient/power scatter rather than a
full mission mix; chronic labels are age-correlated by nature; and one acute episode per
engine at most.
